\section{Lower bound}\label{sec:lower-bound}\label{sec:lower}
Throughout this section $k\ge2$, $T>0$, and $K=\lfloor T\log N\rfloor$, unless a statement allows arbitrary $K$. As above, $P_N^K$ denotes the law of the inverse shuffled permutation and $\nuUnif$ the uniform law on $\mathfrak S_N$. Inversion preserves total variation from uniform, so the same conclusions apply to the chronological forward shuffle. We construct events whose probabilities tend to one under $P_N^K$ and to zero under $\nuUnif$. Every assertion of uniform convergence below means that, after the displayed fixed parameters and an accuracy $\varepsilon>0$ are chosen, one threshold $N_0$ works for every admissible cut array and deterministic family in the statement.
\subsection{Increasing sequences}\subsubsection{Prefix counts}

For a word \(w=(w_1,\ldots,w_m)\) of length \(m\le K\), let
\[
 \lambda_w=\prod_{j=1}^m p^{(j)}_{w_j},\qquad
 t_w=\sum_{\substack{v\in[k]_0^m\\v<_{\rm lex}w}}\lambda_v,
 \qquad C(w)=\#\{i:s_i|_m=w\}.
\]
Here \(s_i|_m\) is the length-\(m\) prefix of the sorted row \(s_i\).
Define the zero-based cuts
\[
 L_w=\#\{i:s_i|_m<_{\rm lex}w\},\qquad
 U_w=\#\{i:s_i|_m\le_{\rm lex}w\}.
\]
Thus the actual ranks occupied by this prefix are
\(I(B_w)=(L_w,U_w]\cap\mathbb Z\), and \(U_w-L_w=C(w)\).
Its deterministic probability interval is
\(J_w=[t_w,t_w+\lambda_w)\). These definitions also apply to the empty
word, whose mass is one and whose count is \(N\).

% Lean source: formalization/lower_bound_v2/RiffleLowerV2/Thinning/Process.lean :: RiffleLowerV2.Thinning.successiveThinning; RiffleLowerV2.Thinning.thinningTails
% Lean source: formalization/lower_bound_v2/RiffleLowerV2/Thinning/Compat/HypergeometricConcentration.lean :: RiffleLowerV2.Thinning.Compat.HypergeomParams.expect_pow_le_all
% Lean source: formalization/lower_bound_v2/RiffleLowerV2/Thinning/Main.lean :: RiffleLowerV2.Thinning.residualPrefixTransform; RiffleLowerV2.Thinning.residualPrefixTails
\begin{lemma}[Successive thinning transforms and tails]
\label{lem:successive-thinning-tail}
Let \(M\ge1\) and let \(P_1,\ldots,P_m\) be independent uniform
subsets of a fixed set of \(M\) labels, with deterministic sizes
\(m_1,\ldots,m_m\). Put
\[
 X=|P_1\cap\cdots\cap P_m|,\qquad
 q_j=m_j/M,\qquad Q=\prod_jq_j,\qquad \mu=MQ.
\]
For an empty intersection use \(X=M\) and \(Q=1\). For every \(s>0\),
\begin{equation}
\label{eq:thinning-transform}
 \E s^X\le(1-Q+Qs)^M\le \exp(\mu(s-1)).
\end{equation}
Consequently, for \(t>0\),
\begin{equation}
\label{eq:thinning-bernstein}
 \Pp(|X-\mu|\ge t)\le
 2\exp\left(-\frac{t^2}{4(\mu+t)}\right),
\end{equation}
and, since $(1-\log2)/2>0$,
\begin{equation}
\label{eq:thinning-lower-tail}
 \Pp(X\le\mu/2)\le e^{-\frac{1-\log2}{2}\mu},
 \qquad \Pp(X<2)\le2e^{-\mu/2}.
\end{equation}
These statements apply to the count of any prescribed word in the shuffle
matrix. They also hold conditionally after any admissible fixed prescriptions
on \(b<N\) labelled rows: delete those rows, replace \(M\) by \(N-b\),
and subtract the prescribed occurrences of each digit from its column count.
\end{lemma}

\begin{proof}
First let $V$ be the intersection of a fixed $x$-element set with a
uniform $m$-element subset of an $M$-element population. Counting prescribed
marked labels gives
\[
 \E\binom Vr=\binom xr\frac{(m)_r}{(M)_r}.
\]
For $s=1+u\ge1$, expand $(1+u)^V$ and use
$(m)_r/(M)_r\le(m/M)^r$. All terms are nonnegative, and hence
\begin{equation}\label{eq:one-step-thinning-transform}
 \E s^V\le(1-m/M+(m/M)s)^x.
\end{equation}
For $0<s\le1$, interchange the roles of the marked set and the sample.
The complement count $x-V$ is hypergeometric with sample size $x$ and
$M-m$ marked positions. Apply the preceding bound to $s^{-1}\ge1$ and
multiply by $s^x$. This gives exactly the right side of
\eqref{eq:one-step-thinning-transform}. The cases $x=0$ and $m=0$ are
included. Here $(n)_r=n(n-1)\cdots(n-r+1)$, with $(n)_0=1$.
Let \(X_j=|P_1\cap\cdots\cap P_j|\), with \(X_0=M\).
Conditionally on the first \(j-1\) subsets,
\(X_j\sim\operatorname{Hyp}(M,X_{j-1},m_j)\).
Starting from \(z_m=s\), set
\(z_{j-1}=1-q_j+q_jz_j\). Every \(z_j\) is positive, and
\(z_0=1-Q+Qs\). Iterating \eqref{eq:one-step-thinning-transform}
gives \(\E s^X\le z_0^M\). The inequality \(1+v\le e^v\)
gives the second bound in \eqref{eq:thinning-transform}.

For \(0\le h\le1\), both
\(e^h-1-h\le h^2\) and \(e^{-h}-1+h\le h^2\).
Substitution of \(s=e^{\pm h}\) into
\eqref{eq:thinning-transform} therefore yields
\(\E e^{\pm h(X-\mu)}\le e^{\mu h^2}\).
Markov's inequality with \(h=t/(2(\mu+t))\) bounds each tail by
\(\exp[-t^2/(4(\mu+t))]\), proving
\eqref{eq:thinning-bernstein}.
Taking \(s=1/2\) in \eqref{eq:thinning-transform} and applying
Markov's inequality at the thresholds \(2^{-\mu/2}\) and \(1/2\)
proves the two bounds in \eqref{eq:thinning-lower-tail}.

In the shuffle matrix each target-digit pile is a uniform fixed-size
subset of the row labels, independently across columns. After prescribed
symbols have been fixed on selected rows, exchangeability makes each
residual column uniform with its residual counts. The conditioning event
factors over columns, so these residual columns remain independent.
The same argument therefore applies with exactly the residual parameters
specified in the statement.
\end{proof}

% Lean source: formalization/lower_bound_v2/RiffleLowerV2/Thinning/Prefixes.lean :: RiffleLowerV2.Thinning.allPrefixConcentration
% Lean source: formalization/lower_bound_v2/RiffleLowerV2/Thinning/UniformError.lean :: RiffleLowerV2.Thinning.prefixFailureUniform
\begin{lemma}[Uniform prefix-count concentration]
\label{cor:iterated-hypergeom}
For every fixed $\varepsilon>0$, the probability that the following inequality fails for any prefix is at most
\[
 2(K+1)k^K\exp(-N^\varepsilon/8)=o(1).
\]
Thus, with probability $1-o(1)$,
simultaneously for every word \(w\) of length at most \(K\),
\begin{equation}
\label{eq:all-prefix-counts}
 |C(w)-N\lambda_w|\le
 N^\varepsilon(\sqrt{N\lambda_w}+1).
\end{equation}
The convergence is uniform over all deterministic column counts.
\end{lemma}

\begin{proof}
Set \(A=N^\varepsilon\), \(\mu=N\lambda_w\), and
\(t=A(\sqrt\mu+1)\). For \(A\ge1\),
\(t^2/(\mu+t)\ge A/2\).
Lemma~\ref{lem:successive-thinning-tail} bounds the failure probability
for this word by \(2e^{-N^\varepsilon/8}\).
There are at most \((K+1)k^{K}\le (T\log N+1)N^{T\log k}\)
words of the indicated lengths. A union bound proves
\eqref{eq:all-prefix-counts}, including words of zero mass, whose counts
are identically zero.
\end{proof}

\subsubsection{The largest word}

% Lean source: formalization/lower_bound_v2/RiffleLowerV2/EmpiricalLower/LargestFinite.lean :: RiffleLowerV2.EmpiricalLower.finiteLargestWordTV
% Lean source: formalization/lower_bound_v2/RiffleLowerV2/EmpiricalLower/LargestAsymptotic.lean :: RiffleLowerV2.EmpiricalLower.largestMean_uniform
% Lean source: formalization/lower_bound_v2/RiffleLowerV2/Thinning/MaximalWord.lean :: RiffleLowerV2.Thinning.maximalMean_eq_rpow_beta
\begin{proposition}[Finite largest-word bound]
\label{prop:largest-prefix}
Let
\(\mu_*=N\prod_{j=1}^{K}\max_\ell p^{(j)}_\ell\).
For every integer \(2\le m\le N\) with \(m\le\mu_*/2\),
\begin{equation}
\label{eq:finite-largest-word-tv}
 \dTV(P_N^{K},\nuUnif)
 \ge1-e^{-\frac{1-\log2}{2}\mu_*}-\frac{N-m+1}{m!}.
\end{equation}
The exact identity $\mu_*=N^{\beta^{(N)}(S_N)}$ holds, where $S_N=K/\log N$. In particular, if \(\mu_*\ge N^b\) for a fixed \(b>0\), the
total-variation distance tends to one, uniformly over these column counts.
\end{proposition}

\begin{proof}
Choose a maximizing digit in every column, breaking ties by the smallest
digit, and let \(w_*\) be the resulting full word. Its count has mean
\(\mu_*\). By \eqref{eq:thinning-lower-tail} it is at least \(m\)
with probability at least \(1-e^{-\frac{1-\log2}{2}\mu_*}\).
Equal full words form one component of the shuffle graph. By
the component-sorting representation, the inverse shuffled permutation
therefore has an increasing run of length \(m\) with at least this
probability. Under a uniform permutation each interval of \(m\) consecutive
positions is increasing with probability \(1/m!\). A union bound over
its \(N-m+1\) possible starting positions proves
\eqref{eq:finite-largest-word-tv}. For the last assertion take
\(m=\lfloor N^{b/2}\rfloor\); the hypotheses force \(b\le1\), and
for all sufficiently large \(N\) this choice satisfies the required
inequalities and \((N-m+1)/m!\to0\).
\end{proof}
\subsection{Cold spots}

For $H\subseteq[N]$, let $E(H)$ be the path edges with both endpoints in $H$, and let $\partial H$ be its boundary in the infinite integer path, including the edges from $0$ to $1$ and from $N$ to $N+1$ when present. For a permutation \(\sigma\in\mathfrak S_N\) and a set \(H\subseteq[N]\),
write
\[
  X_H(\sigma)=\sum_{(i,i+1)\in E(H)}
                 \mathbf 1_{\{\sigma(i)<\sigma(i+1)\}}.
\]
If \(G\) is a subgraph of the path on \([N]\), list its nonempty
components from left to right as \(C_1,\ldots,C_q\).  The permutation
\(\pi_G\) is obtained by sorting the values of \(\pi\) increasingly
on each \(C_j\), and we put \(X_{G,H}(\pi)=X_H(\pi_G)\).
All probability and expectation statements with subscript \(\pi\) below
refer to a uniform permutation, with \(G\) and \(H\) fixed.

% Lean source: formalization/lower_bound_v2/RiffleLowerV2/Ascent/RelativeOrder/Law.lean :: RiffleLowerV2.Ascent.relativeOrderIndependence
% Lean source: formalization/lower_bound_v2/RiffleLowerV2/Ascent/RelativeOrder/Boundary.lean :: RiffleLowerV2.Ascent.boundaryAscentProbability
\begin{lemma}[Independent relative orders]
\label{lem:relative-order-independence}
Let \(A_1,\ldots,A_m\) be pairwise disjoint sets of positions in
\([N]\).  Under a uniform permutation \(\pi\), the relative orders of
the values on these sets are mutually independent and uniform.  In
particular, if \(A\) and \(B\) are disjoint nonempty sets of sizes
\(a\) and \(b\), then
\[
  \Pp_\pi\{\max\pi(A)<\min\pi(B)\}
  =\binom{a+b}{a}^{-1}.
\]
\end{lemma}

\begin{proof}
Put \(d_j=|A_j|\) and \(d_0=N-\sum_{j=1}^m d_j\).  Fix one
relative order on each \(A_j\), using increasing position order to
enumerate its entries.  There are
\(N!/(d_0!\prod_jd_j!)\) ways to assign disjoint value sets of the
prescribed sizes to these position sets and their complement.  Each
assignment admits exactly one ordering on each \(A_j\) with the
prescribed relative order and \(d_0!\) orderings on the complement.
Thus precisely \(N!/\prod_jd_j!\) permutations realize the prescribed
orders.  Dividing by \(N!\) proves the product uniform law.  On
\(A\cup B\), exactly \(a!b!\) of the \((a+b)!\) relative orders
put all values on \(A\) before those on \(B\), proving the last claim.
\end{proof}

% Lean source: formalization/lower_bound_v2/RiffleLowerV2/Ascent/Concentration.lean :: RiffleLowerV2.Ascent.componentConcentration; RiffleLowerV2.Ascent.componentConcentrationTwoSided
\begin{lemma}[Component-sorting concentration]
\label{lem:permutation-bd}
Let \(G\) be a deterministic path subgraph, let \(H\subseteq[N]\),
and let \(M>0\) satisfy \(|H|\le M\).  For every \(t>0\),
\begin{equation}
\label{eq:component-ascent-one-sided}
  \Pp_\pi\{X_{G,H}-\E_\pi X_{G,H}\ge t\},\quad
  \Pp_\pi\{X_{G,H}-\E_\pi X_{G,H}\le-t\}
  \le \exp(-t^2/M).
\end{equation}
Consequently,
\[
  \Pp_\pi\{|X_{G,H}-\E_\pi X_{G,H}|\ge t\}
  \le 2\exp(-t^2/M).
\]
No condition on \(\partial H\) is needed for these estimates.
\end{lemma}

\begin{proof}
Every edge of \(E(G)\cap E(H)\) is an ascent after component
sorting.  Let \(J\subseteq\{1,\ldots,q-1\}\) consist of the
indices for which the edge joining \(C_j\) to \(C_{j+1}\) belongs
to \(E(H)\), and put \(B=|J|\).  Then
\[
  X_{G,H}=|E(G)\cap E(H)|+\sum_{j\in J}Y_j,
  \qquad
  Y_j=\mathbf 1_{\{\max\pi(C_j)<\min\pi(C_{j+1})\}}.
\]
If \(B=0\), this variable is constant and both probabilities in
\eqref{eq:component-ascent-one-sided} vanish.  Assume \(B>0\).
Separate \(J\) into its even and odd component indices.  Within
either part, the sets \(C_j\cup C_{j+1}\) are disjoint, so
Lemma~\ref{lem:relative-order-independence} makes the corresponding
Bernoulli variables mutually independent.  Their means need not equal
\(1/2\).

For completeness, every Bernoulli variable \(Y\) satisfies
\[
  \E\exp\{\lambda(Y-\E Y)\}\le\exp(\lambda^2/8)
  \qquad(\lambda\in\mathbb R).
\]
Indeed the logarithm of the left side has value and first derivative
zero at the origin, and its second derivative is the variance of a
Bernoulli variable under exponential tilting, hence is at most \(1/4\).
Taylor's theorem gives the asserted bound, including the constant
Bernoulli cases.  Let \(S_0,S_1\) be the centered sums over the two
parity classes.  Cauchy--Schwarz and the independence within each class
give
\[
  \E_\pi e^{\lambda(S_0+S_1)}
  \le
  \bigl(\E_\pi e^{2\lambda S_0}
        \E_\pi e^{2\lambda S_1}\bigr)^{1/2}
  \le \exp(\lambda^2B/4).
\]
Markov's inequality, with \(\lambda=2t/B\), bounds either one-sided
tail by \(\exp(-t^2/B)\).  Finally,
\(B\le|E(H)|\le|H|\le M\), which proves
\eqref{eq:component-ascent-one-sided}.  Adding its two estimates proves
the two-sided bound.
\end{proof}

% Lean source: formalization/lower_bound_v2/RiffleLowerV2/Ascent/Expectation.lean :: RiffleLowerV2.Ascent.componentExpectation
\begin{lemma}[Finite component expectation comparison]
\label{lem:component-ascent-expectation}
For every deterministic path subgraph \(G\) and set \(H\subseteq[N]\),
put \(e=|E(H)|\), \(r=|E(G)\cap E(H)|\), and
\(d=|\partial H|\), with the boundary taken in the integer path.
Then
\begin{equation}
\label{eq:component-ascent-expectation}
  \E_\pi X_{G,H}\ge \frac e2+\frac r6-d.
\end{equation}
\end{lemma}

\begin{proof}
Call a component \(C\) of \(G\) partial if it meets both \(H\)
and its complement in \([N]\).  Every partial component contains an
edge of \(\partial H\), and different components contain different
such edges.  Thus there are at most \(d\) partial components.
Call an edge of \(E(H)\setminus E(G)\) bad if at least one of its
two incident components is partial.  A path component is incident to
at most two intercomponent edges, so there are at most \(2d\) bad
edges.  Each has ascent probability at least zero and therefore loses
at most \(1/2\) against the baseline \(1/2\).

The two components incident to any remaining intercomponent edge are
wholly contained in \(H\).  If their sizes are \(a,b\ge1\), its
ascent probability is \(\binom{a+b}{a}^{-1}\), by
Lemma~\ref{lem:relative-order-independence}, and
\begin{equation}
\label{eq:component-boundary-deficit}
  \frac12-\binom{a+b}{a}^{-1}
  \le \frac{(a-1)+(b-1)}6.
\end{equation}
For \(a+b\ge5\), the right side is at least \(1/2\).
For \(a+b=2\) both sides are zero; for \(a+b=3\) both sides are
\(1/6\); and for \(a+b=4\) the left side is either \(1/4\)
or \(1/3\), while the right side is \(1/3\).
Each component wholly contained in \(H\) is charged at most twice
in \eqref{eq:component-boundary-deficit}.  Its \(|C|-1\) internal
edges belong to \(E(G)\cap E(H)\).  The total deficit from these
remaining intercomponent edges is therefore at most \(r/3\).
All \(r\) internal graph edges, including those in partial components,
are retained and contribute excess \(r/2\).  Combining the internal
excess, the deficit at most \(r/3\), and the bad-edge deficit at most
\(d\), gives \eqref{eq:component-ascent-expectation}.
\end{proof}

% Lean source: formalization/lower_bound_v2/RiffleLowerV2/Ascent/Concentration.lean :: RiffleLowerV2.Ascent.uniformScan
% Lean source: formalization/lower_bound_v2/RiffleLowerV2/Ascent/Asymptotic.lean :: RiffleLowerV2.Ascent.uniform_null_scan_eventually
\begin{lemma}[Uniform translate-scan ascent bound]
\label{lem:uniform-ascent}
Let \(H_0\subset\mathbb Z\) be finite with \(M=|H_0|>0\),
and put \(H_q=(q+H_0)\cap[N]\) for \(0\le q\le N\).
For every \(s>0\), the event
\[
  \mathcal A(s)=
  \{\sigma:\exists q\in\{0,\ldots,N\},\
          X_{H_q}(\sigma)>|E(H_q)|/2+s\}
\]
satisfies
\begin{equation}
\label{eq:uniform-ascent-scan-finite}
  \nuUnif(\mathcal A(s))\le (N+1)\exp(-s^2/M).
\end{equation}
In particular, if \(M\ge N^a\) for a fixed \(a>0\), then for
every fixed \(\delta>0\),
\[
  \nuUnif\bigl(\mathcal A(M^{1/2+\delta/2})\bigr)\longrightarrow0.
\]
\end{lemma}

\begin{proof}
A comparison between two distinct positions of a uniform permutation
has mean \(1/2\), so \(\E_\pi X_{H_q}=|E(H_q)|/2\).
Apply Lemma~\ref{lem:permutation-bd} to the edgeless graph and
\(H_q\), for which \(|H_q|\le M\), and take a union bound over
the \(N+1\) choices of \(q\).  This proves
\eqref{eq:uniform-ascent-scan-finite}.  With
\(s=M^{1/2+\delta/2}\), its right side is at most
\((N+1)\exp(-N^{a\delta})\), which tends to zero.
\end{proof}

% Lean source: formalization/lower_bound_v2/RiffleLowerV2/Ascent/Test.lean :: RiffleLowerV2.Ascent.componentComparison
% Lean source: formalization/lower_bound_v2/RiffleLowerV2/Ascent/Quantitative.lean :: RiffleLowerV2.Ascent.quantitative_component_failure
\begin{lemma}[Lalley component comparison]
\label{lem:lalley-comparison}
Let \(M>0\), \(|H|\le M\), and \(G\) be a deterministic
path subgraph.  Write \(e=|E(H)|\),
\(r=|E(G)\cap E(H)|\), and \(d=|\partial H|\).
If \(s>0\) and \(\Delta=r/6-d-s>0\), then
\begin{equation}
\label{eq:lalley-comparison-finite}
  \Pp_\pi\{X_{G,H}\le e/2+s\}
  \le\exp(-\Delta^2/M).
\end{equation}
Consequently, for fixed \(D\ge0\), \(\delta>0\), and any such
family with \(M\to\infty\), \(d\le D\sqrt M\), and
\(r\ge M^{1/2+\delta}\),
\[
  \Pp_\pi\{X_{G,H}>e/2+M^{1/2+\delta/2}\}\longrightarrow1
\]
uniformly over \(N,G,H\).
\end{lemma}

\begin{proof}
Lemma~\ref{lem:component-ascent-expectation} gives
\(\E_\pi X_{G,H}-(e/2+s)\ge\Delta\).
The event in \eqref{eq:lalley-comparison-finite} is therefore
contained in the lower-tail event from
\eqref{eq:component-ascent-one-sided}, with \(t=\Delta\).
For the last assertion take \(s=M^{1/2+\delta/2}\).
For all sufficiently large \(M\), depending only on \(D,\delta\),
\[
  D+M^{\delta/2}\le M^\delta/12,
  \qquad
  \Delta\ge M^{1/2+\delta}/12.
\]
Thus the failure probability is at most
\(\exp(-M^{2\delta}/144)\), uniformly as claimed.
\end{proof}

% Lean source: formalization/lower_bound_v2/RiffleLowerV2/Ascent/Test.lean :: RiffleLowerV2.Ascent.conditionalLalleyTest; RiffleLowerV2.Ascent.conditionalLalleyTendsto
\begin{lemma}[Lalley cold-spot test]
\label{lem:lalley}
Let \(H_0\subset\mathbb Z\) be deterministic and finite, put
\(M=|H_0|\), and let \(H_q=(q+H_0)\cap[N]\).
Suppose that, for constants \(a,\delta>0\) and \(D\ge0\)
independent of \(N\),
\[
  M\ge N^a,\qquad |\partial H_q|\le D\sqrt M
  \quad(0\le q\le N),
\]
and that the shuffle graph satisfies
\[
  \Pp\{\exists q\in\{0,\ldots,N\}:\
         |E(G)\cap E(H_q)|\ge M^{1/2+\delta}\}\longrightarrow1.
\]
Then \(\dTV(P_N^{K},\nuUnif)\longrightarrow1\).
\end{lemma}

\begin{proof}
Use the observable event
\(\mathcal A=\mathcal A(M^{1/2+\delta/2})\) from
Lemma~\ref{lem:uniform-ascent}.  That lemma gives
\(\nuUnif(\mathcal A)\to0\).
Use the representation in the component-sorting representation
to sample the shuffle graph and an independent uniform permutation
\(\pi\), and then form \(\pi_G\).
On the event that a translate has at least
\(M^{1/2+\delta}\) internal graph edges, choose its smallest index
\(q\).  This choice depends only on \(G\).
Conditional on \(G\), Lemma~\ref{lem:lalley-comparison} bounds the
probability that this translate fails the ascent test by
\(\exp(-M^{2\delta}/144)\), for all sufficiently large \(N\).
The bound is uniform in \(G\) and the chosen translate.  Hence
\[
  P_N^K(\mathcal A)
  \ge
  \Pp\{\exists q:\ |E(G)\cap E(H_q)|\ge M^{1/2+\delta}\}
     -\exp(-M^{2\delta}/144)
  \longrightarrow1.
\]
Finally,
\(\dTV(P_N^{K},\nuUnif)
\ge P_N^K(\mathcal A)-\nuUnif(\mathcal A)\to1\).
\end{proof}

\subsection{Construction of the cold spots}
\subsubsection{Empirical costs}

Write $L=\log N$ and $S_N=K/L$.  On
$((j-1)/L,j/L]$, $1\le j\le K$, use the actual profile $p^{(j)}$.
For $0\le r\le s\le S_N$, define
\[
 (A^{(N)},B_\alpha^{(N)},J_\alpha^{(N)},H_\alpha^{(N)})(r,s)
 =\int_r^s(a_{p_N(t)},\psi_{p_N(t)}(\alpha),
                 I_{p_N(t)}(\alpha),H_{p_N(t)}(\alpha))\,dt,
\]
where $p_N(t)=p^{(j)}$ on the displayed column interval.
Here
\[
 p^{[\alpha]}_\ell=\frac{p_\ell^\alpha}{\sum_a p_a^\alpha},\qquad
 I_p(\alpha)=\sum_{\ell:p_\ell>0}p^{[\alpha]}_\ell\log(1/p_\ell),\qquad
 H_p(\alpha)=-\sum_{\ell:p_\ell>0}p^{[\alpha]}_\ell\log p^{[\alpha]}_\ell.
\]
These are integrals of bounded step functions.  In particular, at grid
endpoints they are the normalized finite sums of the actual column costs.
Set
\[
 \beta^{(N)}(s)=1-A^{(N)}(0,s),\qquad
 F_\alpha^{(N)}(s,u)=\frac{\alpha+3}{2}\beta^{(N)}(s)
                 -B_\alpha^{(N)}(s,u)-2B_2^{(N)}(u,S_N).
\]
Set $\Lambda_1^{(N)}=\sup_{0\le s\le u\le S_N}
\inf_{3\le\alpha\le4}F_\alpha^{(N)}(s,u)$.
No distributional assumption on the order of the columns is made here.

For $\alpha\in[2,4]$, the one-column functions satisfy
\begin{equation}\label{eq:repair-global-one-column-bounds}
 0\le a_p\le\log k,\quad
 0\le I_p(\alpha)\le\log k,\quad
 0\le H_p(\alpha)\le I_p(\alpha),\quad
 0\le\psi_p(\alpha)\le(\alpha-1)\log k.
\end{equation}
To verify these assertions, apply the finite Gibbs inequality to
$q=p^{[\alpha]}$ and the weights $p_\ell^\zeta$, restricting all sums to
the positive coordinates. It gives
\[
 H_p(\alpha)\le\zeta I_p(\alpha)-\psi_p(\zeta),\qquad
 H_p(\alpha)=\alpha I_p(\alpha)-\psi_p(\alpha).
\]
Apply the inequality twice, with $\alpha$ and $\zeta$ interchanged. For
$1\le\alpha<\zeta$ their sum gives
$(\zeta-\alpha)(I_p(\zeta)-I_p(\alpha))\le0$.
Thus $I_p$ is decreasing, and $I_p(1)=H(p)\le\log k$; the latter
entropy bound follows from the same Gibbs inequality with the uniform
probability vector. Taking $\zeta=1$ gives $H_p(\alpha)\le I_p(\alpha)$.
Taking $q=p$ in Gibbs gives
$\psi_p(\alpha)\le(\alpha-1)H(p)\le(\alpha-1)\log k$.
Nonnegativity follows from $p_\ell\le1$, and $a_p\le I_p(\alpha)$
follows by averaging $\log(1/p_\ell)\ge-\log\max_a p_a$.
All finite sums are continuous in $\alpha$ on $[2,4]$ and have positive
power-sum denominators. Zero coordinates have zero tilted weight.
Finally, $\sum p_\ell^2\ge(\max p_\ell)^2$ and
$\sum p_\ell^4\le(\sum p_\ell^2)^2$ give
\begin{equation}\label{eq:repair-two-comparisons}
 \psi_p(2)\le2a_p\le2I_p(\alpha),\qquad
 \psi_p(4)\ge2\psi_p(2).
\end{equation}Thus every empirical endpoint displacement by $r$ changes an integrated
$I$ or $H$ by at most $r\log k$, uniformly in the cuts and $N$.


\subsubsection{A finite cold witness}

% Lean source: formalization/lower_bound_v2/RiffleLowerV2/Witness/Main.lean :: RiffleLowerV2.Witness.coldWitness
% Lean source: formalization/lower_bound_v2/RiffleLowerV2/Witness/Continuous.lean :: RiffleLowerV2.Witness.continuousWitness
% Lean source: formalization/lower_bound_v2/RiffleLowerV2/Witness/Finite.lean :: RiffleLowerV2.Witness.roundWitness
\begin{lemma}[Finite cold witness]\label{lem:cold-witness}
Fix $\eta,\rho,\chi_g,\gamma>0$ and assume the local good-column condition on $[0,T]$ with these last three parameters.
Suppose $\Lambda_1^{(N)}\ge\eta$ and
$\beta^{(N)}(S_N)<\min\{\eta/100000,1/1000\}$.
Choose a bridge length $2\rho$ with
$2\rho\log k\le\eta/100000$, where the local good-column hypothesis holds
on every interval of length $\rho$ contained in $[0,T]$. For all sufficiently large $N$ there
are grid endpoints $s=m_s/L$, $u=m_u/L$, a tilt $\theta\in[3,4]$, and an
integer bridge length $b_N=\lceil 2\rho L\rceil$ such that
$0\le m_s\le m_u\le K-b_N$, the bridge columns
$m_u<j\le m_u+b_N$ contain two good columns, and, writing
$b=\beta^{(N)}(s)$, $J=J_\theta^{(N)}(s,u)$,
$H=H_\theta^{(N)}(s,u)$, $J_t=J_2^{(N)}(u,S_N)$,
$H_t=H_2^{(N)}(u,S_N)$,
\begin{gather}
 F_\theta^{(N)}(s,u)\ge\eta/2,\qquad
 J\le b/2-3\eta/1000,\label{eq:repair-witness-upper}\\
 \theta=3\quad\hbox{or}\quad J\ge b/2-\eta/200,
 \label{eq:repair-witness-lower}\\
 h:=b-J+H\ge\eta/1000,\qquad H\le b-J-\eta/1000,
 \label{eq:repair-witness-geometry}\\
 \Gamma:=2b-2J-2J_t+H+H_t,\qquad
 \Gamma-h/2\ge\eta/1000,\label{eq:repair-witness-collision}\\
 b-J-J_t\le\beta^{(N)}(S_N)<\min\{\eta/100000,1/1000\}.
 \label{eq:repair-witness-mean}
\end{gather}
The constants are uniform in the finite cut sequence.
\end{lemma}

\begin{proof}
Choose $s_0\le u_0$ with
$F_\alpha^{(N)}(s_0,u_0)\ge7\eta/8$ for every $\alpha\in[3,4]$.
Write $b_0=\beta^{(N)}(s_0)$.  Nonnegativity of the $B$ terms gives
$b_0\ge\eta/4$, so $b_0/2-\eta/250>0$.
If $J_3^{(N)}(s_0,u_0)\le b_0/2$, choose $\theta=3$.
Otherwise, if $J_4^{(N)}(s_0,u_0)\ge b_0/2$, choose $\theta=4$.
In the remaining case the continuity and monotonicity in $\alpha$ give
$\theta\in(3,4)$ with $J_\theta^{(N)}(s_0,u_0)=b_0/2$.
This selection does not interchange a supremum and an infimum.

If $J_\theta^{(N)}(s_0,u_0)\le b_0/2-\eta/250$, leave $u_1=u_0$;
this case can occur only for $\theta=3$.  Otherwise choose
$u_1\in[s_0,u_0]$ with
$J_\theta^{(N)}(s_0,u_1)=b_0/2-\eta/250$.
Existence follows from continuity and the value zero at $s_0$.
In the $\theta=3$ or interior cases the decrease of $J$ is at most
$\eta/250$.  By \eqref{eq:repair-two-comparisons},
\[
 F_\theta^{(N)}(s_0,u_1)-F_\theta^{(N)}(s_0,u_0)
 =B_\theta^{(N)}(u_1,u_0)-2B_2^{(N)}(u_1,u_0)
 \ge-4J_\theta^{(N)}(u_1,u_0)\ge-2\eta/125.
\]
For $\theta=4$ the same difference is nonnegative, even if the
backward displacement is large.  Hence in every case
$F_\theta^{(N)}(s_0,u_1)\ge7\eta/8-2\eta/125$.

Set
\[
 m_s=\lfloor s_0L\rfloor,\qquad
 m_u=\min\{\lfloor u_1L\rfloor,K-b_N\}.
\]
These integers are ordered for large $N$.  To check the only nontrivial
case, observe that
\[
 (S_N-s_0)\log k\ge A^{(N)}(s_0,S_N)
 =b_0-\beta^{(N)}(S_N)>\eta/4-\min\{\eta/100000,1/1000\}.
\]
Consequently $s_0<S_N-2\rho-2/L$ for large $N$, and clipping $m_u$ to
$K-b_N$ cannot move it below $m_s$.
The endpoint errors satisfy $0\le s_0-s<1/L$ and
$0\le u_1-u\le 2\rho+2/L$.
Equation~\eqref{eq:repair-global-one-column-bounds} bounds the resulting
change of $F$ by $10\rho\log k+20\log k/L$ and that of
$J-b/2$ by $2\rho\log k+4\log k/L$.
These errors are smaller than $\eta/10000$ for large $N$.
The retained slack proves \eqref{eq:repair-witness-upper} and
\eqref{eq:repair-witness-lower}.
For the two good columns, use the interior interval
$[u+\rho/2,u+3\rho/2]$. It lies within the bridge and within $[0,T]$.
Every column time in this interval is strictly above $u$ and at most
$u+b_N/L$, so it belongs to the bridge. The local good-column condition
provides at least $\gamma\rho L>3$ good columns for large $N$; choose two
of them in increasing order.
The exact identities $H=\theta J-B_\theta^{(N)}(s,u)$ and
$H_t=2J_t-B_2^{(N)}(u,S_N)$ imply
\[
 2\Gamma-h=F_\theta^{(N)}(s,u)+(\theta-3)(J-b/2)
 \ge\eta/2-\eta/200.
\]
This is greater than $\eta/500$.
Also $H\le J$ and $J\le b/2-3\eta/1000$ give
$H\le b-J-3\eta/500$ and $h\ge b-J\ge b/2+3\eta/1000$.
Finally $J\ge A^{(N)}(s,u)$ and $J_t\ge A^{(N)}(u,S_N)$ give
\eqref{eq:repair-witness-mean}.
\end{proof}

\subsubsection{Finite type classes}

The construction below uses only the actual columns.  In particular it
needs no temporal approximation for the tilted observables $I$ and $H$.
Fix a $\xi$-regular simplex partition with regularity constant $C\ge0$ and representatives
$p_R$, where $0<\xi<1/(4k)$.  There are $m=m(\xi)$ cells, independent of $N$.
On a chosen set $I$ of column positions, let $I_R$ be the positions in cell
$R$ and $n_R=|I_R|$.  The positions retain their chronological order.
For $\alpha\in[2,4]$ put
\[
 G_R=\{\ell:p_R(\ell)>2\xi\},\qquad
 q_R(\ell)=\frac{p_R(\ell)^\alpha\mathbf1_{G_R}(\ell)}
                 {\sum_{a\in G_R}p_R(a)^\alpha}.
\]
Choose integers $m_{R,\ell}\ge0$ supported on $G_R$, with
$\sum_\ell m_{R,\ell}=n_R$ and
$|m_{R,\ell}-n_Rq_R(\ell)|\le1$; floor the counts and distribute the
remaining units deterministically.  Let $\mathcal C^0_I$ be the words on
$I$ with exactly these counts in each cell.  This class is nonempty.
For every word write $\lambda_x=\prod_{j\in I}p^{(j)}(x_j)$.
Partition $\mathcal C^0_I$ into the integer mass bins
\[
 \{x:j\le-\log\lambda_x<j+1\},\qquad j\in\mathbb Z_{\ge0}.
\]
Select a bin of maximal cardinality, breaking ties by the smallest index, and
call it $\mathcal C_I$.  If its index is $j_I$, define
\[
 J_I^*=j_I/L,\qquad H_I^*=L^{-1}\log|\mathcal C_I|.
\]
These quantities may depend on $N$; that dependence is intentional.

% Lean source: formalization/lower_bound_v2/RiffleLowerV2/Types/Main.lean :: RiffleLowerV2.Types.roundedType
% Lean source: formalization/lower_bound_v2/RiffleLowerV2/Types/Truncation.lean :: RiffleLowerV2.Types.restrictedCost_truncation_error
\begin{lemma}[Finite type estimates]\label{lem:type-estimates}
For $|I|\le TL$ and $\alpha\in[2,4]$, uniformly over partitions with the same fixed number $m$ of cells, the cuts, the chosen
position set, and the selected word,
\begin{gather}
 e^{-1}N^{-J_I^*}<\lambda_x\le N^{-J_I^*},\qquad
 |\mathcal C_I|=N^{H_I^*},\label{eq:repair-bin-exact}\\
 \left|J_I^*-L^{-1}\sum_{j\in I}I_{p^{(j)}}(\alpha)\right|
 +\left|H_I^*-L^{-1}\sum_{j\in I}H_{p^{(j)}}(\alpha)\right|
 \le O_{k,C}(T\xi\log(1/\xi))
       +O_{k,m,\xi,T}\!\left(\frac{\log\log N+1}{\log N}\right).
 \label{eq:repair-types-error}
\end{gather}
Every digit occurring in a selected word has actual column probability
at least $\xi$.  Here the first implied constant depends only on $k,C$, and the second only on $k,m,\xi,T$; these constants and the threshold in $N$ precede the partition, cut array, position set, and tilt. The empty position set is allowed.
\end{lemma}

\begin{proof}
Let $r=p_R$ and let $p$ belong to the same cell. On retained coordinates,
$p_\ell\ge\xi$ and $|\log p_\ell-\log r_\ell|\le C\xi$; on omitted
coordinates, $p_\ell\le3\xi$. The power-sum denominator is at least
$k^{-4}$. The omitted power mass is at most $9k\xi^2$. Moreover,
$p_\ell^\alpha\log(1/p_\ell)\le p_\ell$ for $\alpha\ge2$,
by $\log x\le x-1$, so the omitted cost numerator is at most $3k\xi$.
After division by the full power sum these bounds become $9k^5\xi^2$
and $3k^5\xi$. The retained cost is at most $\log(1/\xi)$. Splitting
the full cost into its retained and omitted parts therefore bounds the
truncation error by $18k^5\xi\log(1/\xi)$.

On retained coordinates, the normalized $\alpha$-power weights at $p$
and $r$ have ratios between $e^{-8C\xi}$ and $e^{8C\xi}$. Comparing
these weights and their digit costs gives
\[
 \left|\sum_\ell q_R(\ell)\log(1/r_\ell)-I_p(\alpha)\right|
 \le O_{k,C}(\xi\log(1/\xi)).
\]
The identity expressing tilted entropy as tilt times cost plus the
logarithm of the power sum, and the same retained/omitted decomposition,
give the corresponding entropy estimate. These bounds are uniform on
$0<\xi<1/(4k)$; this interval also ensures that the retained set is
nonempty, since one coordinate of the representative is at least $1/k$.
Put $\bar J=L^{-1}\sum_R n_R\sum_\ell q_R(\ell)\log(1/p_R(\ell))$
and $\bar H=L^{-1}\sum_R n_RH(q_R)$.
The rounding of counts changes the logarithmic product at representatives
by at most $km\log(1/(2\xi))$.
Replacing representative coordinates by actual coordinates changes it by
at most $C\xi|I|$.  Consequently every $x\in\mathcal C^0_I$ satisfies
\[
 |{-L^{-1}\log\lambda_x}-\bar J|
 \le CT\xi+km\log(1/(2\xi))/L.
\]
The multinomial formula gives
$|\mathcal C^0_I|=\prod_Rn_R!/\prod_\ell m_{R,\ell}!$.
The elementary factorial estimate
$\log(n!)=n\log n-n+O(\log(n+1)+1)$ and the entropy continuity bound
\[
 |n_RH((m_{R,\ell}/n_R)_\ell)-n_RH(q_R)|
 \le O_k(\log(n_R+1)+1)
\]
show the asserted cardinality estimate: if $f(t)=t\log t$, with $f(0)=0$,
then $|f(x)-f(y)|\le\log(n+1)+2$ for $x,y\in[0,n]$ with
$|x-y|\le1$. This follows by applying $\log t\le t-1$ to the ratio of the two positive arguments and using $|t\log t|\le1/e$ on $[0,1]$ when an argument is small. Apply it to
$x=m_{R,\ell}$ and $y=n_Rq_R(\ell)$ and sum over the coordinates;
the empty-cell contribution is zero. Consequently,
$|L^{-1}\log|\mathcal C^0_I|-\bar H|
 =O_{k,m,T}((\log\log N+1)/L)$.
Since every allowed actual digit is at least $\xi$, there are at most
$2+TL\log(1/\xi)$ nonempty integer mass bins.
Keeping a largest bin therefore changes the logarithmic cardinality by
$O_{\xi,T}(\log\log N+1)$ only.  Its index divided by $L$ lies within
$1/L$ of the logarithmic mass of each selected word.
Combining these estimates proves both assertions.
\end{proof}

\subsubsection{Cold sets and bridges}

For a grid endpoint $s=m_s/\log N$, let $a_s$ be the length-$m_s$ word
which uses a maximizing digit in every column, with ties broken by the
smallest digit. For any word on a specified consecutive set of columns,
$\lambda_x$ denotes its product probability on those columns.
For $x$ on $(m_s,m_u]$, put
\[
 t_x^{\rm rel}=\sum_{\substack{x'\in[k]_0^{m_u-m_s}\\x'<_{\rm lex}x}}
                \lambda_{x'},\qquad
 J_x^{\rm rel}=[t_x^{\rm rel},t_x^{\rm rel}+\lambda_x).
\]

% Lean source: formalization/lower_bound_v2/RiffleLowerV2/Localization/Geometry.lean :: RiffleLowerV2.Localization.deterministicColdGeometry
% Lean source: formalization/lower_bound_v2/RiffleLowerV2/Localization/ColdSet.lean :: RiffleLowerV2.Localization.coldRelativeSet_estimates
\begin{lemma}[Finite cold geometry and mass bounds]
\label{lem:cold-geometry}
Fix $\eta,\rho,\chi_g>0$, subject to the
bridge budget $2\rho\log k\le\min\{\eta/100000,1/1000\}/100$. There exist $0<\xi<1/(4k)$ and a $\xi$-regular finite partition with regularity constant $1$, fixed before $N$. For every sufficiently large $N$, every cut array, and every choice of $s,u,\theta,b_N$ satisfying all the conclusions of Lemma~\ref{lem:cold-witness}, there are deterministic prefix
and tail families $\mathcal P_N,\mathcal T_N$, a bridge word $v_b$ with
reserve words $v_b^-,v_b^+$, and a deterministic set $H_{{\rm rel},N}$,
as constructed in the proof. The classes $\mathcal P_N,\mathcal T_N$
are obtained from Lemma~\ref{lem:type-estimates} on the column sets
$(m_s,m_u]$ with tilt $\theta$ and $(m_u+b_N,K]$ with tilt $2$,
respectively; denote their mass-bin exponents by $J_p^*,H_p^*$ and
$J_t^*,H_t^*$. Set
\[
 H_{{\rm rel},N}=\mathbb Z\cap
    \bigcup_{x\in\mathcal P_N}N\lambda_{a_s}J_x^{\rm rel},\qquad
 M_N=|H_{{\rm rel},N}|,\qquad H_q=(q+H_{{\rm rel},N})\cap[N].
\]
Define the real exponents
\[
 \ell_p=b-J_p^*,\quad h_N=\ell_p+H_p^*,\quad
 \ell_b=-\log_N\lambda_{v_b},\quad
 z_N=b-J_p^*-J_t^*-\ell_b,\quad
 g_N=2z_N+H_p^*+H_t^*.
\]
For all sufficiently large $N$,
\begin{gather}
 \ell_p\ge b/2+\eta/500,\qquad H_p^*\le\ell_p-\eta/1000,\qquad
 h_N\ge\eta/2000,\label{eq:cold-effective-geometry}\\
 e^{-2}N^{h_N}\le M_N\le2N^{h_N},\qquad
 N^{\eta/3000}\le M_N\le N,\qquad
 |\partial H_q|\le4\sqrt{M_N}\quad(0\le q\le N),
 \label{eq:cold-geometric-conclusions}\\
 g_N-h_N/2\ge\eta/2000,\qquad z_N\le2\min\{\eta/100000,1/1000\},
 \label{eq:cold-effective-collision}
\end{gather}
The full-word family $\mathcal Y_N=\{a_sxv_bz:x\in\mathcal P_N,z\in\mathcal T_N\}$
satisfies, uniformly in its words,
\begin{equation}
\label{eq:cold-full-word-masses}
 |\mathcal Y_N|=N^{H_p^*+H_t^*},\qquad
 e^{-2}N^{z_N}<N\lambda_y\le N^{z_N}.
\end{equation}
The reserves satisfy $v_b^-<_{\rm lex}v_b<_{\rm lex}v_b^+$ and
\begin{equation}
\label{eq:cold-reserve-means}
 \min_{x\in\mathcal P_N,\,\diamond\in\{-,+\}}
       N\lambda_{a_sxv_b^\diamond}\ge N^{b/2+\eta/1000}.
\end{equation}
Every digit selected by a full word or a reserve has column probability
at least $\min\{\xi,\chi_g,1/k\}>0$. Both type approximations have error at most $\min\{\eta/100000,1/1000\}/1000$. The bridge words differ only at their distinguished columns as described below, and each has mass at least $\chi_g^2k^{-b_N}$. Their central bridge exponent satisfies
\[
 0\le\ell_b\le2\rho\log k+
 \frac{\log k+2\log(1/\chi_g)}{\log N}.
\]
All constants and vanishing errors are uniform
over the finite cut sequences under consideration.
\end{lemma}

\begin{proof}

Use the finite witness, and apply Lemma~\ref{lem:type-estimates} to
$I=(m_s,m_u]$ with $\alpha=\theta$.
Call the resulting class $\mathcal P_N$, with exponents $J_p^*,H_p^*$.
The ambient maximal prefix $a_s$ has exactly
$N\lambda_{a_s}=N^b$ because $s$ is a grid endpoint.
Set
\[
 \ell_p=b-J_p^*,\qquad h_N=\ell_p+H_p^*.
\]
Choose $\xi$, a regular partition with constant $1$, and then $N$ so that each type error in \eqref{eq:repair-types-error} is at most $\min\{\eta/100000,1/1000\}/1000$. The type estimates choose these parameters uniformly before the cut array and the two selected column intervals.
Then
\begin{equation}\label{eq:repair-finite-geometric-buffers}
 \ell_p\ge b/2+\eta/500,\qquad
 H_p^*\le\ell_p-\eta/1000,\qquad h_N\ge\eta/2000.
\end{equation}
For each $x\in\mathcal P_N$ use its exact relative probability interval
$J_x^{\rm rel}$, and put
\[
 H_{{\rm rel},N}=\mathbb Z\cap\bigcup_{x\in\mathcal P_N}
                         N\lambda_{a_s}J_x^{\rm rel}.
\]
The intervals are disjoint and each has length between
$e^{-1}N^{\ell_p}$ and $N^{\ell_p}$.  Their integer counts differ from their
lengths by at most one. Since $N^{\ell_p}\ge e^2$ for sufficiently large $N$, summing the interval counts gives exactly $e^{-2}N^{h_N}\le M_N\le2N^{h_N}$. Therefore, uniformly in $N$,
\[
 |H_{{\rm rel},N}|=N^{h_N+o(1)}\ge N^{\eta/3000},\qquad
 |\partial H_{{\rm rel},N}|\le2N^{H_p^*}
                         =o(|H_{{\rm rel},N}|^{1/2}).
\]
This does not assert that every probability interval is a distinct connected
component: adjacent selected intervals may merge, which only reduces the
boundary.  Translations clipped to $[N]$ have no larger cardinality and
have at most two additional boundary edges. The gap $H_p^*\le\ell_p-\eta/1000$ gives $H_p^*\le h_N/2-\eta/2000$. Hence, for all sufficiently large $N$, $2N^{H_p^*}+2\le4e^{-1}N^{h_N/2}\le4\sqrt{M_N}$, proving the stated coefficient $4$.

Choose two good bridge columns $g_L<g_R$ and define the word $v_b$ as follows.
At $g_L$ use the larger of two good digits; at $g_R$ use the smaller.
Use maximal digits at all other bridge positions.
For the reserve words $v_b^-,v_b^+$, change respectively the digit at
$g_L$ to the smaller and the digit at $g_R$ to the larger, keeping all
other digits equal to the central bridge word.
Then $v_b^-<_{\rm lex}v_b<v_b^+$, and each bridge word has probability at least
$\chi_g^2k^{-b_N}$, where $\chi_g$ is the fixed good-column threshold.
Writing $\ell_b=-L^{-1}\log\lambda_{v_b}$, we obtain
\[
 0\le\ell_b\le 2\rho\log k+(\log k+2\log(1/\chi_g))/L.
\]
The same estimate applies to both reserve words.  In particular every
reserve prefix $a_sxv_b^\pm$ has mean at least
$N^{b/2+\eta/1000}$ for large $N$, after the bridge budget is imposed.
This is the deterministic input for the later concentration/localization
argument; it makes no assertion of independence between reserve counts.

Let $\mathcal R=(m_u+b_N,K]$ and apply
Lemma~\ref{lem:type-estimates} there with $\alpha=2$.
Write $\mathcal T_N,J_t^*,H_t^*$ for the selected tail class and exponents.
The bridge deletion changes the empirical $J_2,H_2$ by at most
$(2\rho+1/L)\log k$ each.
The full selected family is
\[
 \mathcal Y_N=\{a_sxv_bz:x\in\mathcal P_N, z\in\mathcal T_N\}.
\]
For $y\in\mathcal Y_N$ put $\mu_y=N\lambda_y$ and define
\[
 z_N=b-J_p^*-J_t^*-\ell_b,\qquad
 g_N=2z_N+H_p^*+H_t^*.
\]
The mass bins give the exact uniform bounds
\[
 e^{-2}N^{z_N}<\mu_y\le N^{z_N},\qquad
 |\mathcal Y_N|=N^{H_p^*+H_t^*},\qquad
 \sum_{y\in\mathcal Y_N}\mu_y^2=N^{g_N+o(1)}.
\]
Write
\[
 \delta_{p,J}=J_p^*-J,\quad \delta_{p,H}=H_p^*-H,\quad
 \delta_{t,J}=J_t^*-(J_t-d_J),\quad
 \delta_{t,H}=H_t^*-(H_t-d_H),
\]
where $d_J=J_2^{(N)}(u,u+b_N/L)$ and
$d_H=H_2^{(N)}(u,u+b_N/L)$ are the deleted bridge integrals.
Each $\delta$ has absolute value at most $\min\{\eta/100000,1/1000\}/1000$ and
$0\le d_J,d_H\le d:=(2\rho+1/L)\log k$.
For all sufficiently large $N$, our budgets give
$d,\ell_b\le\min\{\eta/100000,1/1000\}/50$. Direct subtraction now yields
\begin{align*}
 |h_N-h|&\le\frac1{500}\min\{\eta/100000,1/1000\},\\
 g_N-h_N/2&\ge\Gamma-h/2
   -\frac7{1000}\min\{\eta/100000,1/1000\}-d-2\ell_b
   \ge\eta/2000.
\end{align*}
and
\begin{align*}
 z_N&\le b-J-J_t+\frac1{500}\min\{\eta/100000,1/1000\}+d\\
    &<\frac{501}{500}\min\{\eta/100000,1/1000\}+d\\
    &\le2\min\{\eta/100000,1/1000\}.
\end{align*}
These prove \eqref{eq:cold-effective-collision}. The reserve mass lower
bound proved above gives \eqref{eq:cold-reserve-means}.
All selected digits have probability at least
$c=\min\{\xi,\chi_g,1/k\}>0$. Finally
$H_{{\rm rel},N}\subset\mathbb Z\cap[0,N\lambda_{a_s})\subset[0,N)\cap\mathbb Z$,
so $M_N\le N$. The bounds already established give
\eqref{eq:cold-geometric-conclusions}, including its uniform clipped-boundary
estimate, and complete the proof.
\end{proof}

\subsubsection{Localization}

% Lean source: formalization/lower_bound_v2/RiffleLowerV2/Localization/Endpoints.lean :: RiffleLowerV2.Localization.relativeEndpoint
\begin{lemma}[Finite relative endpoint estimate]
\label{lem:finite-relative-endpoints}
Let \(a\) be a deterministic prefix of length \(l\), let
\(l\le m\le K\), and put \(\mu_a=N\lambda_a\).
For a suffix \(x\) on columns \(l+1,\ldots,m\), let
\(\lambda_x\) be its product mass and \(t_x\) the total product
mass of the length-\((m-l)\) suffixes strictly smaller than \(x\).
Define \(L_x^{\mathrm{rel}},U_x^{\mathrm{rel}}\) by counting rows
with prefix \(a\) whose length-\((m-l)\) suffixes are respectively
smaller than, or at most, \(x\).
If \(A>0\) and
\[
 |C(v)-N\lambda_v|\le A(\sqrt{N\lambda_v}+1)
 \quad\hbox{for every prefix \(v\) of length at most \(K\)},
\]
then simultaneously for all such \(a,x\),
\begin{equation}
\label{eq:relative-endpoint-finite-error}
\begin{split}
 &|L_x^{\mathrm{rel}}-\mu_at_x|
  +|U_x^{\mathrm{rel}}-\mu_a(t_x+\lambda_x)|\\
 &\hspace{6mm}\le
 A\bigl[(2\sqrt{kK}+1)\sqrt{\mu_a}+2kK+1\bigr].
\end{split}
\end{equation}
\end{lemma}

\begin{proof}
Partition the preceding and following suffixes by the first column
where they differ from \(x\).  The resulting exact identities are
\begin{align}
\label{eq:relative-left-side-exact}
 L_x^{\mathrm{rel}}-\mu_at_x
 &=\sum_{r=l+1}^m\sum_{d<x_r}
   \bigl[C(ax_{l+1}\cdots x_{r-1}d)
          -N\lambda_{ax_{l+1}\cdots x_{r-1}d}\bigr],\\
\label{eq:relative-right-side-exact}
 U_x^{\mathrm{rel}}-\mu_a(t_x+\lambda_x)
 &=C(a)-\mu_a
  -\sum_{r=l+1}^m\sum_{d>x_r}
   \bigl[C(ax_{l+1}\cdots x_{r-1}d)
          -N\lambda_{ax_{l+1}\cdots x_{r-1}d}\bigr].
\end{align}
Each side contains at most \(kK\) disjoint subcylinders of the
ambient cylinder.  Their means \(\mu_j\) therefore have sum at
most \(\mu_a\).  Cauchy--Schwarz bounds either side sum by
\[
 A\sum_j(\sqrt{\mu_j}+1)
 \le A\bigl[\sqrt{kK\mu_a}+kK\bigr].
\]
The additional ambient fluctuation in
\eqref{eq:relative-right-side-exact} is at most
\(A(\sqrt{\mu_a}+1)\).  Adding these bounds proves
\eqref{eq:relative-endpoint-finite-error}, including empty suffixes
and zero or small means.
\end{proof}

% Lean source: formalization/lower_bound_v2/RiffleLowerV2/Localization/Events.lean :: RiffleLowerV2.Localization.localization; RiffleLowerV2.Localization.coldLocalization_of_prefix_event
\begin{lemma}[Uniform cold-word localization]
\label{lem:localization}
Use the finite witness, the selected prefix family \(\mathcal P_N\),
the full-word family \(\mathcal Y_N\), and the cold set
\(H_{{\rm rel},N}\).  Put
\[
 H_q=(q+H_{{\rm rel},N})\cap[N],\qquad q_*=L_{a_s},\qquad
 b=\beta^{(N)}(s),\qquad N\lambda_{a_s}=N^b.
\]
With probability \(1-o(1)\), simultaneously for every
\(y\in\mathcal Y_N\), all rows carrying \(y\) have sorted
positions in \(H_{q_*}\).  On this event each distinct full word
that occurs at least twice supplies a distinct edge of
\(E(G)\cap E(H_{q_*})\).  The probability bound is uniform over
the permitted finite witnesses and cut arrays.
\end{lemma}

\begin{proof}
Take \(\varepsilon=\eta/4000\) in
Lemma~\ref{cor:iterated-hypergeom}.  On its simultaneous
prefix-count event, Lemma~\ref{lem:finite-relative-endpoints} bounds
both relative endpoint errors for every \(x\in\mathcal P_N\)
by
\[
 e_N=N^{\eta/4000}
       \bigl[(2\sqrt{kK}+1)N^{b/2}+2kK+1\bigr].
\]
Since \(J_p^*\ge0\), the inequality
\(\ell_p=b-J_p^*\ge b/2+\eta/500\) in
\eqref{eq:repair-finite-geometric-buffers} implies \(b\ge\eta/250\).
As \(K\le T\log N\), uniformly for all these witnesses,
\begin{equation}
\label{eq:repair-endpoint-error-concrete}
 e_N\le N^{b/2+\eta/2000}
 \quad\hbox{for all sufficiently large \(N\)}.
\end{equation}

The reserve words \(v_b^-\) and \(v_b^+\) differ from the central
bridge \(v_b\) at only the first and second distinguished good
column, respectively.  At the distinguished columns their digit
probabilities are at least \(\chi_g\), and all other bridge digits
are maximal.  Hence their bridge products are at least
\(\chi_g^2k^{-b_N}\), where
\(b_N=\lceil 2\rho\log N\rceil\) is the bridge length.  By
\eqref{eq:repair-bin-exact} and
\eqref{eq:repair-finite-geometric-buffers}, every reserve prefix has
mean at least
\[
 \mu_x^\pm:=\E C(a_sxv_b^\pm)
 \ge e^{-1}\chi_g^2 k^{-1}
       N^{\ell_p-2\rho\log k}
 \ge N^{b/2+\eta/1000}
\]
for all sufficiently large \(N\), because
\(2\rho\log k\le\eta/100000\).  The fixed factor involving
\(\chi_g\) is absorbed by the remaining positive exponent gap.

The same simultaneous prefix event already controls every reserve count.
For a reserve mean $\mu\ge N^{b/2+\eta/1000}$, use $b\ge\eta/250$ and,
for sufficiently large $N$, $16\le N^{\eta/400}$ to obtain
\[
 N^{\eta/4000}(\sqrt\mu+1)\le\mu/2.
\]
Indeed $\mu\ge16N^{\eta/2000}$, so $\sqrt\mu\ge4N^{\eta/4000}$;
also $\sqrt\mu\ge1$. The prefix deviation bound gives $C(v)\ge\mu/2$.
The further inequality $4\le N^{\eta/2000}$ and the bound on $e_N$ yield
\begin{equation}\label{eq:repair-reserve-count-dominates}
 C(a_sxv_b^-),\ C(a_sxv_b^+)\ge\tfrac12N^{b/2+\eta/1000}
                           \ge e_N+2
 \qquad(x\in\mathcal P_N).
\end{equation}
Thus one prefix-count event controls both endpoints and all reserves,
with failure probability at most
$2(K+1)k^K\exp(-N^{\eta/4000}/8)$. Its threshold depends only on the
fixed parameters, since it is enough that
$4kT\log N+4\le N^{\eta/4000}$, $4\le N^{\eta/2000}$, and
$16\le N^{\eta/400}$.

Work on this event.For a fixed \(x\), write
\[
 \ell=N\lambda_{a_s}t_x^{\rm rel},\qquad
 r=N\lambda_{a_s}(t_x^{\rm rel}+\lambda_x).
\]
The relative endpoints differ from \(\ell,r\) by at most
\(e_N\).  Every one-based sorted row rank \(i\) carrying a word
\(a_sxv_bz\in\mathcal Y_N\) satisfies
\[
 q_*+L_x^{\mathrm{rel}}+C(a_sxv_b^-)+1
 \le i\le
 q_*+U_x^{\mathrm{rel}}-C(a_sxv_b^+),
\]
because \(v_b^-<_{\rm lex}v_b<_{\rm lex}v_b^+\).
Equation~\eqref{eq:repair-reserve-count-dominates} therefore gives
\(q_*+\ell<i<q_*+r\).  Thus
\[
 i\in q_*+(\mathbb Z\cap[\ell,r))\subseteq H_{q_*}.
\]
The two spare units in the reserve bound resolve the one-based-rank
and half-open-interval conventions.  No union bound over tail words
\(z\) and no estimate on the random translation \(q_*\) are needed.

Copies of a complete word are consecutive after lexicographic sorting.
If there are at least two, choose the edge between its first two copies.
The localization places both ranks in \(H_{q_*}\), and this pair
is an edge of \(G\).  Distinct words occupy disjoint constant-word
blocks, so their chosen edges are distinct.
\end{proof}
\subsubsection{Collision counts}

For a full word $w$, write $N_w=C(w)$ for its occupancy and
$\mu_w=N\lambda_w$ for its mean.

% Lean source: formalization/lower_bound_v2/RiffleLowerV2/Factorial/Exact.lean :: RiffleLowerV2.Factorial.exactJointFactorial
% Lean source: formalization/lower_bound_v2/RiffleLowerV2/Factorial/RelativeTargets.lean :: RiffleLowerV2.Factorial.relativeJointFactorial
\begin{lemma}[Fixed-column factorial moments]
\label{lem:fixed-column-factorial}
For arbitrary $K$, distinct full words $w_1,\ldots,w_q$, and nonnegative
integers $r_i$ with $R=\sum_i r_i\le N$, let $m_{j,a}=\sum_{i:(w_i)_j=a}r_i$.
Then the exact identity is
\[
 \E\prod_i(N_{w_i})_{r_i}
   =(N)_R\prod_{j=1}^K\frac{\prod_a(n_a^{(j)})_{m_{j,a}}}{(N)_R}.
\]
In addition, suppose $K\le C\log N$ and every pile requested by the distinct full words
$w_1,\ldots,w_q$ has size at least $cN$, for fixed $c,C>0$.
For nonnegative integers $r_i$ with $R=\sum_i r_i\le N^{1/5}$,
\begin{equation}
\label{eq:exact-fixed-column-factorials}
 \E\prod_{i=1}^q(N_{w_i})_{r_i}
 =\prod_{i=1}^q\mu_{w_i}^{r_i}
    \left(1+O_{c,C}\left(\frac{(K+1)R^2}{N}\right)\right).
\end{equation}
The bound is uniform in the words and the column counts.
\end{lemma}

\begin{proof}
Expanding the falling factorials chooses ordered labelled rows for each
word. As the words are distinct, only disjoint choices contribute. There
are $(N)_R$ such choices. Let $m_{j,a}$ be the number of their requests for
digit $a$ in column $j$; thus $\sum_a m_{j,a}=R$.
Uniform sampling of each fixed-count column gives the exact formula
\[
 \E\prod_i(N_{w_i})_{r_i}
 =(N)_R\prod_{j=1}^K
                \frac{\prod_a(n_a^{(j)})_{m_{j,a}}}{(N)_R}.
\]
For $n\ge cN$ and $0\le v\le R\le N^{1/5}$,
\[
 \log\frac{(n)_v}{n^v}
 =\sum_{h=0}^{v-1}\log(1-h/n)=O_c(v^2/N),
\]
uniformly for sufficiently large $N$. Apply this also with $n=N$,
and compare the exact formula with
\[
 \prod_i\mu_{w_i}^{r_i}
   =N^R\prod_{j,a}(n_a^{(j)}/N)^{m_{j,a}}.
\]
Since $\sum_a m_{j,a}^2\le R^2$, the logarithm of their ratio is
$O_c((K+1)R^2/N)=o(1)$. Exponentiating proves
\eqref{eq:exact-fixed-column-factorials}. The case $R=0$ is equality.
\end{proof}

% Lean source: formalization/lower_bound_v2/RiffleLowerV2/Factorial/Central.lean :: RiffleLowerV2.Factorial.centralComparison; RiffleLowerV2.Factorial.centralComparisonOne
\begin{lemma}[Finite factorial comparison with a controlled remainder]
\label{lem:central-window}
Let $X,Y$ be nonnegative integer-valued random variables on a finite probability space, let $\mu,\nu>0$,
and let $P_\mu,P_\nu$ be independent Poisson variables. Suppose that $m\ge2$
is even and, for $r,s\ge0$ with $r+s\le2m$,
\[
 \left|\E(X)_r(Y)_s-\mu^r\nu^s\right|
 \le\epsilon\mu^r\nu^s,\qquad 0\le\epsilon\le1.
\]
Set $f(t)=1-(1+t)e^{-t}$ and $d_m(t)=(m-1)t^m/m!$. Then
\begin{align*}
 \left|\Pp(X\ge2,Y\ge2)-f(\mu)f(\nu)\right|
 &\le \epsilon(1+\mu)(1+\nu)e^{\mu+\nu}\\
 &\quad+3\bigl(d_m(\mu)+d_m(\nu)+d_m(\mu)d_m(\nu)\bigr).
\end{align*}
If only the one-variable moments through order $m$ are assumed, then
\[
 \left|\Pp(X\ge2)-f(\mu)\right|
 \le\epsilon(1+\mu)e^\mu+3d_m(\mu).
\]
\end{lemma}

\begin{proof}
Define the even Bonferroni polynomial and its last term by
\[
 U_m(n)=\sum_{r=2}^{m}(-1)^r(r-1)\binom nr,
 \qquad D_m(n)=(m-1)\binom nm.
\]
The finite binomial identities give the pointwise inequalities
\begin{equation}
\label{eq:collision-bonferroni-sandwich}
 0\le U_m(n)-\mathbf 1_{\{n\ge2\}}\le D_m(n)
 \qquad(n\ge0).
\end{equation}
For completeness, put $S_j(n)=\sum_{r=2}^{j}(-1)^r(r-1)\binom nr$.
For $n>j\ge2$,
\[
 S_j(n)-1=(-1)^j
 \left(n\binom{n-2}{j-1}-\binom{n-1}{j}\right),
\]
and the expression in parentheses equals
$((j-1)n+1)\binom{n-2}{j-1}/j\ge0$.
For $2\le n\le j$, $S_j(n)=1$; for $n=0,1$, it is zero.
Thus $S_{m-1}(n)\le\mathbf 1_{\{n\ge2\}}\le S_m(n)$; the case $m=2$
uses $S_1=0$. Their difference is $D_m(n)$, proving
\eqref{eq:collision-bonferroni-sandwich}.

Writing $I_X=\mathbf 1_{\{X\ge2\}}$ and similarly for $Y$, the same inequality gives
\[
 0\le U_m(X)U_m(Y)-I_XI_Y
 \le D_m(X)+D_m(Y)+D_m(X)D_m(Y).
\]
Consequently the expected error is at most
$(1+\epsilon)(d_m(\mu)+d_m(\nu)+d_m(\mu)d_m(\nu))$.
The independent Poisson pair has the same bound with $\epsilon=0$.
The difference between the expectations of the two polynomial products is
at most
\[
 \epsilon\left(\sum_{r=2}^{m}\frac{(r-1)\mu^r}{r!}\right)
 \left(\sum_{s=2}^{m}\frac{(s-1)\nu^s}{s!}\right)
 \le\epsilon(1+\mu)(1+\nu)e^{\mu+\nu}.
\]
Adding the three bounds proves the assertion. The one-variable assertion
follows from the same sandwich. All expectations here are of polynomials of
total degree at most $2m$: in particular their values on $\{X+Y>2m\}$ have
been controlled, rather than discarded using only the probability of that event.
\end{proof}

% Lean source: formalization/lower_bound_v2/RiffleLowerV2/Occupancy/Main.lean :: RiffleLowerV2.Occupancy.uniformOccupancy; RiffleLowerV2.Occupancy.collisionChebyshev
\begin{lemma}[Uniform fixed-column collision probabilities]
\label{lem:two-category-poisson}
Fix $c>0$ and $C>0$. For each $N$, let $K\le C\log N$, and take $K$
independent uniformly permuted fixed-count columns on the same $N$ labelled
rows. Let $\mathcal W$ be any deterministic family of distinct full words
such that every pile selected by any word in $\mathcal W$ has size at least
$cN$. Put
\[
 \lambda_w=\prod_{t=1}^{K}\frac{n_{w_t}^{(t)}}N,
 \quad \mu_w=N\lambda_w,
 \quad A_w=\{N_w\ge2\},\quad p_w=\Pp(A_w),
 \quad f(t)=1-(1+t)e^{-t}.
\]
Uniformly over all such data and words, as $N\to\infty$,
\begin{align}
 p_w&=f(\mu_w)\bigl(1+O_{c,C}(N^{-1/100})\bigr),
 \label{eq:uniform-collision-marginal}\\
 \Pp(A_w\cap A_{w'})
 &=p_wp_{w'}\bigl(1+O_{c,C}(N^{-1/100})\bigr),\qquad w\ne w'.
 \label{eq:uniform-collision-pair}
\end{align}
In particular $p_w\asymp\min\{\mu_w^2,1\}$. For
$Z=\sum_{w\in\mathcal W}\mathbf 1_{A_w}$,
\begin{equation}
\label{eq:uniform-collision-variance}
 \Var Z\le\E Z+O_{c,C}(N^{-1/100})(\E Z)^2.
\end{equation}
No upper or lower restriction on $\mu_w$ is required.
\end{lemma}

\begin{proof}
All constants below depend only on $c,C$. The exact labelled-row calculation
in Lemma~\ref{lem:fixed-column-factorial} gives, for distinct words $w,w'$
and $r+s\le R\le N^{1/5}$,
\begin{equation}
\label{eq:collision-uniform-factorials}
 \E(N_w)_r(N_{w'})_s
 =\mu_w^r\mu_{w'}^s\left(1+O\left(\frac{(K+1)R^2}{N}\right)\right).
\end{equation}
The same one-variable estimate holds in every model obtained by forcing
$b\le3$ labelled rows to carry a word in $\mathcal W$ and deleting these
rows: the residual population is $N-b$, the columns remain independent
uniform fixed-count columns, and every selected residual pile has size at
least $(c/2)(N-b)$ for all sufficiently large $N$.

For the occupancy $V$ of a fixed word in any of these original or residual
models, let $v=\E V$. Equation~\eqref{eq:thinning-transform} gives
\begin{equation}
\label{eq:collision-empty-singleton-tail}
 \E s^V\le e^{-(1-s)v}\quad(0<s\le1),\qquad
 \Pp(V<2)\le2e^{-v/2}.
\end{equation}
The latter bound is also \eqref{eq:thinning-lower-tail}.

Choose once and for all
\[
 \tau=N^{-1/100},\qquad B=\frac1{10}\log N,
 \qquad R=\lfloor N^{1/25}\rfloor,
 \qquad m=2\lfloor R/4\rfloor.
\]
For large $N$, $m\ge2$ is even and $2m\le R\le N^{1/5}$.
The relative factorial error in \eqref{eq:collision-uniform-factorials}
is bounded by
\[
 \epsilon=O_{c,C}((\log N)N^{-23/25}).
\]
If $0<\mu,\nu\le B$, Lemma~\ref{lem:central-window} therefore gives
\begin{equation}
\label{eq:collision-central-absolute}
 \left|\Pp(N_w\ge2,N_{w'}\ge2)-f(\mu)f(\nu)\right|
 =O_{c,C}(N^{-7/10}).
\end{equation}
Its one-variable version gives the same bound whenever $0<\mu\le2B$,
including the residual models on $N-b$ rows, $b\le3$.
To verify these errors explicitly, the moment-error term is at most
$C(\log N)^3N^{-18/25}=O(N^{-7/10})$ in the two-variable range,
and at most $C(\log N)^2N^{-18/25}$ in the one-variable range $\mu\le2B$.
Moreover,
\[
 d_m(t)\le m\exp(2t-m\log2)\le N^{-1},\qquad 0<t\le2B,
\]
for all sufficiently large $N$. Indeed $m\le N$, $m\log2\ge3\log N$, and $2t\le(2/5)\log N$ suffice for this estimate. The product of two such terms is at most $N^{-2}$. Changing $N$ to $N-b$ changes only the constants.

We first prove \eqref{eq:uniform-collision-marginal}.
If $\mu\le\tau$, the order-two/order-three Bonferroni inequalities are
\begin{equation}
\label{eq:collision-rare-bonferroni}
 \frac12\E(V)_2-\frac13\E(V)_3
 \le\Pp(V\ge2)\le\frac12\E(V)_2.
\end{equation}
The coefficient of $(V)_3$ is $1/3$.
Equation~\eqref{eq:collision-uniform-factorials} thus gives
$p_w=\mu^2(1+O(\tau+(K+1)/N))/2$.
Since $f(\mu)=\mu^2(1+O(\mu))/2$, the desired relative estimate follows,
uniformly even when $\mu$ is arbitrarily small.
If $\tau<\mu\le B$, then $f(\mu)\ge(e^{-1}/2)\min\{\mu^2,1\}\ge(e^{-1}/2)\tau^2$.
The one-variable central bound has relative error
$O(N^{-7/10}\tau^{-2})=O(N^{-17/25})$, again smaller than $O(\tau)$.
If $\mu>B$, \eqref{eq:collision-empty-singleton-tail} gives
$1-p_w\le2N^{-1/20}$, whereas
$1-f(\mu)=(1+\mu)e^{-\mu}\le(1+B)e^{-B}=O((\log N)N^{-1/10})$.
This proves \eqref{eq:uniform-collision-marginal} in all three regimes.
The same conclusion holds uniformly in the residual models, by the verified
population and pile-size bounds above.

This marginal result also provides the conditional stability needed when
one of the two categories is rare. Fix $b=2$ or $3$ distinct labelled rows
and force them to carry $w$. For $w'\ne w$, none of the forced rows counts
as a $w'$ row. The conditional mean of the remaining $w'$ count is
\[
 \nu_b=(N-b)\prod_{t=1}^{K}
 \frac{n_{w'_t}^{(t)}-b\mathbf 1_{\{w'_t=w_t\}}}{N-b},
 \qquad \frac{\nu_b}{\nu}=1+O((K+1)/N),\quad\nu=\mu_{w'}.
\]
All these conditioning events are positive for large $N$ because their
requested piles have size at least $cN$. The function $f$ satisfies
\[
 0\le\frac{tf'(t)}{f(t)}\le2\qquad(t>0),
\]
since $f(t)=\int_0^t u e^{-u}\,du\ge e^{-t}t^2/2$.
Equivalently, $f$ is increasing and $f(t)/t^2$ is decreasing. Thus
$0<x\le y$ implies $f(x)\le f(y)\le(y/x)^2f(x)$. Applying these two
inequalities to the original and residual means gives
$f(\nu_b)/f(\nu)=1+O((K+1)/N)$ uniformly for every $\nu>0$.Applying \eqref{eq:uniform-collision-marginal} to the original and residual
models proves
\begin{equation}
\label{eq:collision-conditional-stability}
 \Pp(A_{w'}\mid\text{the prescribed $b$ rows carry $w$})
 =p_{w'}(1+O(\tau)),\qquad b=2,3.
\end{equation}
The constants are independent of the row labels and both words.

To prove \eqref{eq:uniform-collision-pair}, first suppose
$\mu=\mu_w\le\tau$; no restriction is imposed on $\nu=\mu_{w'}$.
Summing \eqref{eq:collision-conditional-stability} over the ordered labelled
$b$-tuples gives
\[
 \E[(N_w)_b\mathbf 1_{A_{w'}}]
 =\mu^b p_{w'}(1+O(\tau)),\qquad b=2,3.
\]
Multiplying the pointwise Bonferroni sandwich
\eqref{eq:collision-rare-bonferroni} by $\mathbf 1_{A_{w'}}$ therefore yields
\[
 \Pp(A_w\cap A_{w'})=\frac{\mu^2}2p_{w'}(1+O(\tau))
 =p_wp_{w'}(1+O(\tau)).
\]
This includes the case where both means are arbitrarily small.

It remains to consider $\mu,\nu>\tau$. If, say, $\nu>B$, then
\eqref{eq:collision-empty-singleton-tail} implies
\[
 |\Pp(A_w\cap A_{w'})-p_wp_{w'}|
 \le2\Pp(A_{w'}^c)\le4N^{-1/20}.
\]
Since $p_w\ge(e^{-1}/4)\tau^2$ and $p_{w'}\ge1/2$ for large $N$, its relative
error is $O(N^{-3/100})=O(\tau)$.
Finally, if $\tau<\mu,\nu\le B$, use
\eqref{eq:collision-central-absolute} and the already proved marginal
relative estimates. Here $f(\mu)f(\nu)\ge(e^{-2}/4)\tau^4$, so the additional
central relative error is
$O(N^{-7/10}\tau^{-4})=O(N^{-33/50})$.
This proves the pair estimate uniformly in every case.

It follows that
\[
 \sum_{w\ne w'}|\operatorname{Cov}(\mathbf 1_{A_w},\mathbf 1_{A_{w'}})|
 \le C\tau\left(\sum_w p_w\right)^2.
\]
Adding $\sum_w\Var(\mathbf 1_{A_w})\le\E Z$ proves
\eqref{eq:uniform-collision-variance}.
\end{proof}

% Lean source: formalization/lower_bound_v2/RiffleLowerV2/EmpiricalLower/ManyEdges.lean :: RiffleLowerV2.EmpiricalLower.manyEdges
% Lean source: formalization/lower_bound_v2/RiffleLowerV2/EmpiricalLower/CollisionMean.lean :: RiffleLowerV2.EmpiricalLower.cold_collision_mean_lower
\begin{lemma}[Many localized collision edges]
\label{lem:many-edges}
For the finite witness and cold set of Lemma~\ref{lem:cold-geometry},
the following estimate holds:
\[
 \Pp\{|E(G)\cap E(H_{L_{a_s}})|\ge M_N^{1/2+\eta/5000}\}
 \longrightarrow1.
\]
The convergence is uniform over the permitted finite witnesses and cuts.
\end{lemma}

\begin{proof}
Let $A_y=\{N_y\ge2\}$ and $Z=\sum_{y\in\mathcal Y_N}\mathbf 1_{A_y}$.
Every selected pile has size at least $cN$ for a fixed $c>0$, by
Lemma~\ref{lem:cold-geometry}. Hence
Lemma~\ref{lem:two-category-poisson} applies to this whole deterministic
family. Its marginal bound and \eqref{eq:cold-full-word-masses} give
\[
 \Pp(A_y)\ge \frac{e^{-1}}4\min\{\mu_y^2,1\}
 \ge \frac{e^{-5}}4N^{2z_N-4\min\{\eta/100000,1/1000\}},\qquad
 \E Z\ge \frac{e^{-5}}4N^{g_N-4\min\{\eta/100000,1/1000\}},
\]
Here we used the upper bound on $z_N$ in \eqref{eq:cold-effective-collision}.
The other bound there gives
\[
 g_N-h_N/2-4\min\{\eta/100000,1/1000\}
 \ge\eta/2000-\eta/25000>\eta/3000.
\]
Together with \eqref{eq:cold-geometric-conclusions} this yields, for all
sufficiently large $N$,
\[
 \E Z\ge M_N^{1/2}N^{\eta/3000},\qquad \E Z\longrightarrow\infty.
\]
Equation~\eqref{eq:uniform-collision-variance} and Chebyshev's inequality
now imply
\[
 \Pp(Z<\tfrac12\E Z)
 \le\frac4{\E Z}+O(N^{-1/100})=o(1).
\]
In particular $Z\ge M_N^{1/2}N^{\eta/4000}$ with probability $1-o(1)$.
Intersect this event with that of Lemma~\ref{lem:localization}. Its
injective assignment of edges gives
$|E(G)\cap E(H_{L_{a_s}})|\ge Z$. Since $M_N\le N$, any fixed
$0<\delta<\eta/4000$, for example $\delta=\eta/5000$, proves the claim.
\end{proof}






\subsubsection{Completion of the empirical lower bound}

% Lean source: formalization/lower_bound_v2/RiffleLowerV2/EmpiricalLower/Cold.lean :: RiffleLowerV2.EmpiricalLower.coldOnly
% Lean source: formalization/lower_bound_v2/RiffleLowerV2/EmpiricalLower/Cold.lean :: RiffleLowerV2.EmpiricalLower.coldOnly
\begin{proposition}[Cold-spot lower bound]\label{prop:cold-lower}
Fix $\eta,\chi_g,\gamma,\rho>0$ with
\[
 2\rho\log k\le\frac1{100}\min\{\eta/100000,1/1000\}.
\]
For every $\varepsilon>0$, there exists $N_0\ge2$ such that, for every
$N\ge N_0$ and every cut array satisfying the local good-column condition
with parameters $\chi_g,\gamma,\rho$ on $[0,T]$, the inequalities
\[
 \Lambda_1^{(N)}\ge\eta,\qquad
 \beta^{(N)}(S_N)<\min\{\eta/100000,1/1000\}
\]
imply $\dTV(P_N^K,\nuUnif)\ge1-\varepsilon$.
\end{proposition}

\begin{proof}
Lemma~\ref{lem:cold-witness} supplies the finite witness, and
Lemma~\ref{lem:cold-geometry} supplies its deterministic cold set. Its
cardinality is at least $N^{\eta/3000}$, and every clipped translate has
boundary at most four times the square root of that cardinality.
Lemma~\ref{lem:many-edges} gives at least $M_N^{1/2+\eta/5000}$ internal
graph edges on the translate indexed by $L_{a_s}$, with probability tending
to one uniformly. Apply Lemma~\ref{lem:lalley} with $a=\eta/3000$,
$\delta=\eta/5000$, and $D=4$. The thresholds for the geometry, collisions,
localization, and ascent test are all fixed before the cut array. Taking
their maximum gives the asserted $N_0$.
\end{proof}

% Lean source: formalization/lower_bound_v2/RiffleLowerV2/EmpiricalLower/Main.lean :: RiffleLowerV2.EmpiricalLower.empiricalCombined
% Lean source: formalization/lower_bound_v2/RiffleLowerV2/EmpiricalLower/Main.lean :: RiffleLowerV2.EmpiricalLower.empiricalCombined


\begin{proof}[Proof of Theorem~\ref{thm:empirical-lower}]
The first alternative follows from Proposition~\ref{prop:largest-prefix}
and the exact maximal-word identity. In the second alternative, if
$\beta^{(N)}(S_N)\ge\min\{\eta_1/100000,1/1000\}$, apply the same
proposition with this fixed positive exponent. If the reverse strict
inequality holds, apply Proposition~\ref{prop:cold-lower} with
$\eta=\eta_1$. Take the maximum of the three thresholds. When the cold
parameters are not positive or fail the displayed restriction, only the
first alternative can occur, and its threshold alone suffices.
\end{proof}

\subsubsection{Passage to a macroscopic profile}

We now apply the profile transfer proved in Proposition~\ref{prop:profile-transfer}.
Its interval and margin bounds, \eqref{eq:interval-transfer} and
\eqref{eq:margin-transfer}, use the same selected finite partition in the
almost-profile condition and the scalar integrability hypotheses. In particular,
\[
 |\beta^{(N)}(S_N)-\Lambda_0(T)|\le e+\frac{\log k}{\log N},\qquad
 |\Lambda_1^{(N)}-\Lambda_1(T)|\le\frac{13e}{2}+\frac{14\log k}{\log N}.
\]
% Lean source: formalization/lower_bound_v2/RiffleLowerV2/MacroscopicLower/Main.lean :: RiffleLowerV2.MacroscopicLower.macroscopicIntegrable
% Lean source: formalization/lower_bound_v2/RiffleLowerV2/MacroscopicLower/Main.lean :: RiffleLowerV2.MacroscopicLower.macroscopicIntegrable_tendsto
% Lean source: formalization/lower_bound_v2/RiffleLowerV2/MacroscopicLower/Main.lean :: RiffleLowerV2.MacroscopicLower.macroscopicIntegrable; RiffleLowerV2.MacroscopicLower.macroscopicIntegrable_tendsto


\begin{proof}[Proof of Theorem~\ref{thm:macroscopic-lower}]
Choose
\[
 0<e<\min\left\{\frac1{10}\min\{\eta/200000,1/1000\},
                          \frac{\eta}{100}\right\}
\]
and apply Proposition~\ref{prop:profile-transfer} with a mesh upper bound
small enough for the displayed restriction on $\rho_0$. All these choices
precede the profile and the cut array. If
$\beta^{(N)}(S_N)\ge\min\{\eta/200000,1/1000\}$, the largest-word
proposition applies with this fixed positive exponent. Otherwise,
\eqref{eq:margin-transfer} gives $\Lambda_0(T)<\eta$ for
large $N$. The assumed maximum therefore implies $\Lambda_1(T)\ge\eta$.
Equation~\eqref{eq:margin-transfer} then gives
$\Lambda_1^{(N)}\ge\eta/2$ for large $N$. Proposition~\ref{prop:cold-lower}
with margin $\eta/2$ applies in this second case. Take the maximum of the
transfer and the two lower-bound thresholds. The inequalities
$0\le\dTV\le1$ give the asserted convergence for eventually admissible
sequences.
\end{proof}





